%%%%%%%%%%%%%%%%%%%%%%%%%%%%%%%%%%%%%%%%%
% Beamer Presentation
% LaTeX Template
% Version 1.0 (10/11/12)
%
% This template has been downloaded from:
% http://www.LaTeXTemplates.com
%
% License:
% CC BY-NC-SA 3.0 (http://creativecommons.org/licenses/by-nc-sa/3.0/)
%
%%%%%%%%%%%%%%%%%%%%%%%%%%%%%%%%%%%%%%%%%

%----------------------------------------------------------------------------------------
%	PACKAGES AND THEMES
%----------------------------------------------------------------------------------------

\documentclass{beamer}

\mode<presentation> {

% The Beamer class comes with a number of default slide themes
% which change the colors and layouts of slides. Below this is a list
% of all the themes, uncomment each in turn to see what they look like.

%\usetheme{default}
%\usetheme{AnnArbor}
%\usetheme{Antibes}
%\usetheme{Bergen}
%\usetheme{Berkeley}
%\usetheme{Berlin}
%\usetheme{Boadilla}
%\usetheme{CambridgeUS}
%\usetheme{Copenhagen}
%\usetheme{Darmstadt}
%\usetheme{Dresden}
%\usetheme{Frankfurt}
%\usetheme{Goettingen}
%\usetheme{Hannover}
%\usetheme{Ilmenau}
%\usetheme{JuanLesPins}
%\usetheme{Luebeck}
%\usetheme{Madrid}
%\usetheme{Malmoe}
%\usetheme{Marburg}
%\usetheme{Montpellier}
%\usetheme{PaloAlto}
\usetheme{Pittsburgh}
%\usetheme{Rochester}
%\usetheme{Singapore}
%\usetheme{Szeged}
%\usetheme{Warsaw}

% As well as themes, the Beamer class has a number of color themes
% for any slide theme. Uncomment each of these in turn to see how it
% changes the colors of your current slide theme.

%\usecolortheme{albatross}
\usecolortheme{beaver}
%\usecolortheme{beetle}
%\usecolortheme{crane}
%\usecolortheme{dolphin}
%\usecolortheme{dove}
%\usecolortheme{fly}
%\usecolortheme{lily}
%\usecolortheme{orchid}
%\usecolortheme{rose}
%\usecolortheme{seagull}
%\usecolortheme{seahorse}
%\usecolortheme{whale}
%\usecolortheme{wolverine}

%\setbeamertemplate{footline} % To remove the footer line in all slides uncomment this line
%\setbeamertemplate{footline}[page number] % To replace the footer line in all slides with a simple slide count uncomment this line

%\setbeamertemplate{navigation symbols}{} % To remove the navigation symbols from the bottom of all slides uncomment this line
}

\usepackage{graphicx} % Allows including images
\usepackage{booktabs} % Allows the use of \toprule, \midrule and \bottomrule in tables


% \usefonttheme{professionalfonts} % using non standard fonts for beamer
% \usefonttheme{serif} % default family is serif
% \usepackage{fontspec}
% \setmainfont{Calibri} % 这四行命令修改全局字体

\usepackage{ragged2e} % 两端对齐
\justifying\let\raggedright\justifying
%----------------------------------------------------------------------------------------
%	TITLE PAGE
%----------------------------------------------------------------------------------------

\title[Asset Pricing]{\textbf{Asset Pricing}} % The short title appears at the bottom of every slide, the full title is only on the title page

\subtitle{5. Beta}

\author[Cheng Hang]{Cheng Hang} % Your name
\institute[DUFE] % Your institution as it will appear on the bottom of every slide, may be shorthand to save space
{School of Fintech \\
Dongbei University of Finance \& Economics \\ % Your institution for the title page
\medskip
\textbf{chenghangch@qq.com}% Your email address
}
\date{\today} % Date, can be changed to a custom date

	\AtBeginSection[]
{
	\begin{frame}{Content}
		\tableofcontents[sectionstyle=show/shaded,subsectionstyle=show/shaded/hide] %这几个参数我也不知道该如何准确地解释，反正它们最终的效果是突出显示当前章节，而其它章节都进行了淡化处理
		\addtocounter{framenumber}{-1}  %目录页不计算页码
	\end{frame}
}

%----------------------------------------------------------------------------------------
%	Begin
%----------------------------------------------------------------------------------------

\begin{document}

\begin{frame}
\titlepage % Print the title page as the first slide
\end{frame}

%\begin{frame}
%\frametitle{Overview} % Table of contents slide, comment this block out to remove it
%\tableofcontents % Throughout your presentation, if you choose to use \section{} and \subsection{} commands, these will automatically be printed on this slide as an overview of your presentation
%\end{frame}

%----------------------------------------------------------------------------------------
%	PRESENTATION SLIDES
%----------------------------------------------------------------------------------------

\section{Time-Series}
\subsection{Forecast Stock Market Return}

\begin{frame}{Forecast Market Return: In Sample}
The commonly used \textbf{Predictive Regression} is, For simplicity, considering a univariate predictive regression of the period-($t + 1$) stock market return (excess return) $r_{t+1}$ on a single predictor variable $x_t$:
\begin{equation}\label{equ: 预测收益率}
	r_{t+1}=\alpha+\beta x_t+\varepsilon_{t+1} \quad \text { for } t=1, \ldots, T-1
\end{equation}

There are two primary reasons for the widespread use of the simple predictive regression in Equation \textcolor{blue}{(\ref{equ: 预测收益率})}.
\begin{enumerate}
	\item It is straightforward to implement, with the results intuitively understandable in a familiar regression framework.
	\item It should capture some return predictability, even if the true DGP is more complex.
\end{enumerate}
\end{frame}

\begin{frame}{Forecast Market Return}
\begin{itemize}
	\item $R^2$ statistic for the predictive regression in Equation \textcolor{blue}{(\ref{equ: 预测收益率})} is usually quite small.
	\item It is worth noting that the conventional ordinary least squares (OLS) estimator of $\beta$ in Equation \textcolor{blue}{(\ref{equ: 预测收益率})} is generally biased.

		\begin{itemize}
			\item \textcolor{blue}{\href{https://doi.org/10.1016/S0304-405X(99)00041-0}{Stambaugh (1999)}}
			\item \textcolor{blue}{\href{https://doi.org/10.1093/rfs/hhu139}{Kostakis, Magdalinos, and Stamatogiannis (2015)}}
			\item \textcolor{blue}{\href{https://doi.org/10.1016/j.jfineco.2019.08.004}{Huang et al. (2020)}}
		\end{itemize}
\end{itemize}
\end{frame}

\begin{frame}{Out-of-Sample Tests of Return Predictability}
	This in sample regression can be misleading in that using all the data leads to a ``lookahead" bias, because an investor in real time does not have access to all the sample.
	As a result, the traditional in-sample approach cannot be used to make forecasts that mimic the situation of an investor in real time.

	An out-of-sample forecast of $r_{t+1}$ based on the predictor $x_t$ and information available through period $t$ is given by
	\begin{equation}
		\hat{r}_{t+1 \mid t}=\hat{\alpha}_t+\hat{\beta}_t x_t
	\end{equation}
	where $\hat{\alpha}_t$ and $\hat{\beta}_t$ are the OLS estimates of $\beta$ and $\alpha$, respectively, in Equation \textcolor{blue}{(\ref{equ: 预测收益率})} based on data available through $t$.
\end{frame}

\subsection{Test Pricing Factors}

\begin{frame}{Test Pricing Factors}
	The commonly used method is
	\begin{equation}
		R_{i t}=\alpha_i+\beta_i^{\prime} f_t+\varepsilon_{i t}, \quad t=1,2, \cdots, T, \forall i
    \end{equation}
	where $R_{i t}$ is the stock or portfolio $i$\('\)s return, $f_t$ is the factor (return). $\beta_i$ is the exposure corresponding to the factor $f_t$.

	The average of the time-series is
	\begin{equation}
		\mathrm{E}_T\left[R_i\right]=\alpha_i+\beta_i^{\prime} \mathrm{E}_T\left[f_t\right]
	\end{equation}

Note: You need to estimate the model using individual portfolios or stock returns.

\end{frame}

\section{Cross-Section}

\begin{frame}{Cross Section Regression}
The factors in cross section regression can be economic variables, such as inflation, GDP, money supply, and technological shock. So, the first step of cross section regression is time-series estimation:
\begin{equation}\label{cross section reg first}
	R_{i t}=a_i+\beta_i^{\prime} f_t+\varepsilon_{i t}, \quad t=1,2, \cdots, T, \forall i
\end{equation}
where $f_t$ is the pricing factor. $beta_i^{\prime}$ is the factor loading or exposure. The next step is
\begin{equation}\label{cross section reg second}
	\mathrm{E}\left[R_i\right]=\beta_i^{\prime} \lambda+\alpha_i, \quad i=1,2, \cdots, N
\end{equation}
Note:
\begin{enumerate}
	\item In the previous model, the left hand is the expected return (the mean of return within $T$ period), and the right-hand regressor is $\beta$ estimated by time-series regression of model \textcolor{blue}{(\ref{cross section reg first})}.
	\item There is no intercept in model \textcolor{blue}{(\ref{cross section reg second})}. $\alpha_i$ is the pricing error.
	\item $\beta$ is esitmated.
\end{enumerate}

\end{frame}

\subsection{Fama-MacBeth Regression}

\begin{frame}{Fama-MacBeth Regression}
	The initial stage of Fama-MacBeth regression involves acquiring the exposure of each asset on all factors via $N$ time series regression, which is analogous to the first step of cross-sectional regression. The dissimilarity between the Fama-MacBeth regression and the sectional regression test primarily manifests in the second step of sectional regression. At each time point $t$, a Fama-MacBeth regression is conducted where the dependent variable is the expected return (specifically, the return of the period and not the average rate of all periods). During the $t$ time period, the cross-sectional linear regression model for asset excess return and factor return is expressed as follows:

	\begin{equation}
		R_{i t}^e=\gamma_t+\hat{\boldsymbol{\beta}}_i^{\prime} \boldsymbol{\lambda}_t+\alpha_{i t}, \quad i=1,2, \cdots, N
	\end{equation}

\end{frame}

\begin{frame}{Fama-MacBeth Regression}
\begin{equation}
	\begin{aligned}
		\hat{\boldsymbol{\lambda}} & =\frac{1}{T} \sum_{t=1}^T \hat{\boldsymbol{\lambda}}_t \\
		\hat{\alpha}_i & =\frac{1}{T} \sum_{t=1}^T \hat{\alpha}_{i t}
		\end{aligned}
\end{equation}
\end{frame}

\begin{frame}{Main Steps}
	\begin{enumerate}
		\item Summary Statistics
		\begin{enumerate}
			\item Cross-Sectional Summary Statistics
			\item Time-series averages of the cross-sectional summary statistics
		\end{enumerate}
		\item Correlation
		\begin{enumerate}
			\item Cross-sectional Pearson product-moment ($\rho_t$) and Spearman rank ($\rho^S_t$) correlations
			\item Time-series averages of the cross-sectional correlation coefficient
		\end{enumerate}
	\end{enumerate}
\end{frame}

\section{Factors}

\begin{frame}{Factors and Factors Exposures}

	\begin{itemize}
		\item Improve the estimation of factors exposures
		\begin{itemize}
			\item Fama and MacBeth (1973)
		\end{itemize}
		\item How to construct factors?
		\begin{itemize}
			\item Factor mimicking portfolio
			\begin{enumerate}
				\item The portfolio only has exposure to the target factor and zero exposure to other factors;
				\item Of all the portfolios satisfying condition 1, this portfolio has the least idiosyncratic risk.
			\end{enumerate}
		\end{itemize}
	\end{itemize}
\end{frame}

\begin{frame}
\begin{table}[H]
	\small
	\centering
	\resizebox{!}{!}{
		\begin{tabular}{cccc}
			\hline \\[-2ex]
			         & exposure factor A & exposure factor B &  idiosyncratic volatility \\
			\hline \\ [-2ex]
			stock 1  & 1 & 1.2 &  1\%  \\
			stock 2  & 1 & 1 &   1\%  \\
			stock 3  & 2 & -1.2 &  1\% \\
			stock 4  & 1.1 & -1.2 &  5\% \\
			stock 5  & 0.1 & -1.2 &  5\% \\

			\hline \\
			\end{tabular}

}
\end{table}
\end{frame}

\begin{frame}{Factors: Portfolio Sort}
	The primary benefit of utilizing the \textbf{Portfolio Sort} method is its ability to eliminate the prerequisite that the factor exposure is already established.

	This approach presupposes solely that variables and factor exposures exhibit an association. Using $BM$ as an example case, this approach suggests that equities with higher BM exhibit greater susceptibility to value-based factors. Conversely, equities with lowered BM display reduced susceptibility to the value-based factors without any additional implications. Within the framework of this fundamental concept, individuals may construct a factor portfolio based on the level of $BM$, even in the absence of knowledge regarding the exposure of individual stocks value factor.
\end{frame}

\begin{frame}[allowframebreaks]{Factors: Portfolio Sort}
Main steps:
\begin{enumerate}
	\item Sort: The stocks are organized in the cross section based on the value of the sorting variable (BM, in this instance) in descending order.
	\begin{itemize}
		\item Three portfolios: 30th, 70th; Five portfolios: 20th, 40th, 60th, 80th
	\end{itemize}
	\item Group: The construction of a spread portfolio involves the establishment of a long-short hedge portfolio, whereby the highest-ranked group stocks are purchased while simultaneously shorting the lowest-ranked group stocks. (Fund neutrality, value-weighted or equal-weighted)
	\begin{equation}
	r_p = r_{high} - r_{low}
	\end{equation}
	\item Rebalance: Monthly or Yearly (Trading Cost)
\end{enumerate}
Then, we test
\begin{itemize}
	\item Significance: The null hypothesis is
\begin{equation*}
	r_p = 0
\end{equation*}
	\item Monotonicity
\end{itemize}
\end{frame}

\section{China's Stock Market Data}

\begin{frame}{Trading Rules}
	\begin{itemize}
		\item 1996.12.16 $\pm 10\%$ daily price limit (10.4\%), except
		\begin{itemize}
			\item The IPO date
			\item The first trading date after the stock split-structure reform
			\item The first trading date after seasonal offerings
			\item The first trading date after material assets restructuring.
			\item The first re-listing date of de-listed stocks
			\item Science and technology innovation board: $\pm 20\%$ daily price limit after initial 5 days
		\end{itemize}
		\item 1995 ``T+1'' rules
		\item Since April 22, 1998, the Shanghai and Shenzhen Stock Exchange announced that firms with financial abnormality will undergo special treatment with “ST” being added as a prefix to its stock name.
		\item The Split-Share Structure and Its Reform
		\item March 2010, margin trading and short selling (MTSS) program
	\end{itemize}
\end{frame}

\begin{frame}{Common Data Filters}
	\begin{enumerate}
		\item Financial firms
		\item IPO stocks
		\item ST stocks
		\item Trading days
	\end{enumerate}

\end{frame}










\end{document}